% Section content template for individual Educative lessons
% This file contains the actual content for a specific section
% Generated from Educative API section content

% Section: Key Insights and Tips: Designing an Elevator System
% Section ID: 4973261040254976
% Section Slug: key-insights-and-tips-designing-an-elevator-system
% Generated: 2025-12-11T14:08:55.007845

Congratulations on completing the case study on designing an elevator system. To help solidify your understanding, this lesson will cover common design mistakes you should avoid and specific tips for your interview. Afterward , you can test your knowledge with a short quiz and explore related case studies to see these design principles applied in different contexts.

\subsection{Common mistakes}\label{Dpo0jLsNagCPsaDyZPyTG}
\\
Avoiding these common mistakes will help you develop a more efficient and scalable elevator system:

\begin{itemize}
\item
\protect\phantomsection\label{OOuI-Jf82YeuJBt8tzZ9b}
\textbf{Hardcoding the number of floors and elevators:} Designing the system with a fixed number of elevators and floors limits scalability. You should use configurable parameters to accommodate different building sizes.
\item
\protect\phantomsection\label{4hy7cglQKZQA2u2HHQszt}
\textbf{Ignoring the dispatcher algorithm:} Not implementing a clear strategy for which elevator responds to a call can result in inefficient passenger wait times. You should implement a dispatcher algorithm, like selecting the nearest idle elevator or an elevator moving toward the passenger.
\item
\protect\phantomsection\label{ryr803A6_aavo5FF9jCxV}
\textbf{Failing to model doors separately:} Treating the door as a simple attribute of the elevator car oversimplifies its functionality. Create a distinct \texttt{Door} class with its own states (\texttt{Open}, \texttt{Closed}) and methods to manage its operation independently.
\item
\protect\phantomsection\label{QjwhvSO2Rq0ITqFxryVXI}
\textbf{Not providing separate displays:} Assuming the same information is needed inside and outside the elevator is a design flaw. Design separate display functionalities, one for the hall showing direction and arrival, and a more detailed one inside showing the current floor, state, and capacity.
\item
\protect\phantomsection\label{PNtKSNBlwLw5AOvQsha9t}
\textbf{Ignoring concurrency issues:} Assuming only one request happens at a time is a critical oversight in a multi-elevator, multi-floor system. You should use appropriate data structures, like a request queue, to handle concurrent requests from multiple floors and inside multiple elevators.
\end{itemize}

\subsection{Interview tips}\label{GfHryhjT5altWo7oqo-Nh}
\\
Here are some essential tips to help you prepare for an interview where you're asked to design an elevator system:

\begin{center}
\begin{tabular}{|p{0.3\linewidth}|p{0.65\linewidth}|}
\hline
\textbf{Tip} & \textbf{Explanation} \\
\hline
Use enumerations for states. & Clearly define enumerations for ElevatorState (IDLE, UP, DOWN), DoorState (OPEN, CLOSED), and Direction (UP, DOWN). This makes your code cleaner, more readable, and less error-prone than simple strings or integers. \\
\hline
Design for extension using the Open/Closed principle. & Mention how using the Strategy pattern for dispatching allows new algorithms to be added without changing the ElevatorSystem class. This shows you are designing a system that can evolve without modifying its core code. \\
\hline
Start with a bottom-up design approach. & Explain that you will start by designing the smallest components, like the Button and Door, and then compose them into larger components like ElevatorPanel, Floor, and ElevatorCar. This demonstrates a structured and methodical design process. \\
\hline
Address safety and constraints. & Explain how your design handles constraints like maximum capacity (680 kg or 8 persons). Discussing the isOverloaded() check and alarm shows you are considering real-world safety requirements. \\
\hline
Talk about extensibility. & Briefly discuss how your system could be extended, for instance, by adding new types of elevators (e.g., freight) or more sophisticated dispatching strategies. This shows forward-thinking and an understanding of scalable design. \\
\hline
\end{tabular}
\end{center}


\subsection{Test your understanding}\label{w9HMfXs-YjkGp6pHHON6C}
\\
Take a short quiz to evaluate your understanding of the elevator system design.

\subsection*{Designing an Elevator System}
\vspace{0.3cm}
\noindent
\textbf{Question 1:} Which design pattern is most suitable for handling the various elevator states, such as \texttt{IDLE}, \texttt{UP}, and \texttt{DOWN}?
\\[0.2cm]
\begin{itemize}
 \item Factory pattern
 \item Singleton pattern
 \item Observer pattern
 \item \textbf{[CORRECT]} State pattern
\\[0.1cm]
\textit{Explanation:} 
The State pattern is ideal for managing an object's behavior when its internal state changes, and this pattern directly applies to the elevator's different operational states.
\end{itemize}
\vspace{0.3cm}
\noindent
\textbf{Question 2:} What is the primary responsibility of the \texttt{dispatcher()} component in the \texttt{ElevatorSystem}?
\\[0.2cm]
\begin{itemize}
 \item Opening and closing the elevator doors
 \item \textbf{[CORRECT]} Selecting the most appropriate elevator car to service a request
\\[0.1cm]
\textit{Explanation:} 
The dispatcher contains the algorithm to intelligently select the best elevator car to minimize wait times and optimize system efficiency.
 \item Displaying the current floor number
 \item Moving the elevator car between floors
\end{itemize}
\vspace{0.3cm}
\noindent
\textbf{Question 3:} How does the system prevent the ``Up'' button from being used on the topmost floor?
\\[0.2cm]
\begin{itemize}
 \item The \texttt{HallButton} is disabled in the \texttt{ElevatorPanel}.
 \item The \texttt{Dispatcher} ignores requests from the top floor.
 \item The \texttt{ElevatorCar} automatically rejects the request.
 \item \textbf{[CORRECT]} The \texttt{Floor} class has an \texttt{isTopMost()} method that can be used to disable the button.
\\[0.1cm]
\textit{Explanation:} 
The \texttt{Floor} class is aware of its position and includes logic (\texttt{isTopMost()} and \texttt{isBottomMost()}) to control the state of its hall panel buttons.
\end{itemize}
\vspace{0.3cm}
\noindent
\textbf{Question 4:} What kind of relationship do \texttt{HallButton} and \texttt{ElevatorButton} have with the \texttt{Button} class?
\\[0.2cm]
\begin{itemize}
 \item \textbf{[CORRECT]} Inheritance
\\[0.1cm]
\textit{Explanation:} 
\texttt{HallButton} and \texttt{ElevatorButton} are specialized types of a \texttt{Button} and extend its basic properties, which is a classic use of inheritance.
 \item Association
 \item Aggregation
 \item Composition
\end{itemize}
\vspace{0.3cm}
\noindent
\textbf{Question 5:} Which class is responsible for composing the \texttt{HallPanel} and \texttt{Display} on each floor?
\\[0.2cm]
\begin{itemize}
 \item \texttt{Building}
 \item \texttt{ElevatorSystem}
 \item \textbf{[CORRECT]} \texttt{Floor}
\\[0.1cm]
\textit{Explanation:} 
Each \texttt{Floor} is composed of its own set of \texttt{HallPanels} and \texttt{Displays}, as specified in the requirements.
 \item \texttt{ElevatorCar}
\end{itemize}

\subsection{Related case studies}\label{CzKYJwdQAtExgHT9kYvul}
\\
Explore these related case studies to deepen your understanding of the design principles applied in the elevator system:

\begin{itemize}
\item
\protect\phantomsection\label{X5yVy2xFiwHsq2CY8W1tC}
\textbf{Designing a parking lot:} This case study also involves managing a finite set of resources (parking spots vs. elevators) and processing user requests based on availability and proximity.
\item
\protect\phantomsection\label{ebyMM0i_2RIXnoAPQpcir}
\textbf{Designing a vending machine:} This problem requires careful state management (e.g., \texttt{Idle}, \texttt{AcceptingCoins}, \texttt{DispensingItem}) and handling user input through a panel, similar to the elevator's states and panels.
\item
\protect\phantomsection\label{nqQk1DiM5CfeJI7iiy0ci}
\textbf{Designing a traffic control system:} This is a classic system design problem that involves managing the flow of entities (vehicles) through a network of nodes (intersections) using algorithms to optimize throughput and minimize wait times, analogous to the elevator dispatcher.
\item
\protect\phantomsection\label{KO3XLzkuWzOPUvuIXxjad}
\textbf{Designing a job scheduler for a computing cluster:} This problem requires an intelligent dispatcher to assign incoming computing jobs to the most appropriate server based on resource availability and priority, mirroring how an elevator system assigns cars to passenger requests.
\item
\protect\phantomsection\label{e6d0GEirnY8clgDmhHCf3}
\textbf{Designing a ride sharing service (e.g., Uber):} This involves a central system that dispatches the nearest available driver (resource) to a user's request, which is a direct parallel to the core logic of the elevator dispatch system.
\end{itemize}

% End of section content